% SEGMENT OLP-0016-S01
% Part: sets-functions-relations
% Chapter: relations
% Section: orders

\documentclass[../../../include/open-logic-section]{subfiles}

\begin{document}

\olfileid{sfr}{rel}{ord}
\olsection{வரிசைகள்}

\begin{explain}
நாம் செய்யும் பல ஒப்பீடுகளில், ஏதோ ஒரு வகையில் சில பொருள்கள் மற்றவற்றைவிட
“சிறியவை”, அவற்றுக்கு “சமமானவை”, அல்லது அவற்றைவிட “பெரியவை” என்று
விவரிக்கிறோம். இவற்றில் \emph{வரிசைத்} தொடர்புகள் இடம்பெறுகின்றன. ஆனால்
வரிசைத் தொடர்புகள் பல வகைப்படும். எடுத்துக்காட்டாக, எந்த இரு பொருள்களையும்
ஒப்பிட இயல வேண்டும் என்று சில வகைகள் கோருகின்றன; மற்றவை அவ்வாறு
கோருவதில்லை. சில ஒன்றாதலை உள்ளடக்குகின்றன ($\le$ போல); சில அதை
விலக்குகின்றன ($<$ போல). இங்கு ஒரு வகைப்பாடு நமக்குப் பயன்படும்.
\end{explain}

% SEGMENT OLP-0016-S02
\begin{defn}[முன்வரிசை]
தற்சுட்டுப் பண்பையும் கடப்புப் பண்பையும் கொண்ட ஒரு தொடர்பு
\emph{முன்வரிசை} எனப்படும்.
\end{defn}

% SEGMENT OLP-0016-S03
\begin{defn}[பகுதி வரிசை]
எதிர்சமச்சீர்ப் பண்பையும் கொண்ட ஒரு முன்வரிசை \emph{பகுதி வரிசை}
எனப்படும்.
\end{defn}

% SEGMENT OLP-0016-S04
\begin{defn}[நேரியல் வரிசை]\ollabel{def:linearorder}
ஒப்பிடத்தக்க பண்பையும் கொண்ட ஒரு பகுதி வரிசை \emph{முழு வரிசை}
அல்லது \emph{நேரியல் வரிசை} எனப்படும்.
\end{defn}

% SEGMENT OLP-0016-S05
\begin{ex}
ஒவ்வொரு நேரியல் வரிசையும் ஒரு பகுதி வரிசையாகும்; ஒவ்வொரு பகுதி வரிசையும்
ஒரு முன்வரிசையாகும். ஆனால் இவற்றின் மறுதலைகள் உண்மையல்ல.
$A$ மீதான அனைத்துத் தொடர்பு தற்சுட்டுப் பண்பையும் கடப்புப் பண்பையும்
கொண்டதால் அது ஒரு முன்வரிசை. ஆனால் $A$ ஒன்றுக்கு மேற்பட்ட
!!{element}[acc]க் கொண்டிருந்தால், அனைத்துத் தொடர்புக்கு எதிர்சமச்சீர்ப்
பண்பு இல்லை; எனவே அது பகுதி வரிசையல்ல.
\end{ex}

% SEGMENT OLP-0016-S06
\begin{ex}
$\Bin^*$ மீதான \emph{நீளம் மிகாத} தொடர்பான $\preccurlyeq$ ஐக்
கருதுவோம்: $x \preccurlyeq y$ என்பது $\len{x} \le \len{y}$ ஆக
இருக்கும்போது, அப்போதும் அப்போது மட்டுமே, உண்மையாகும். இது ஒரு
முன்வரிசை (தற்சுட்டுப் பண்பும் கடப்புப் பண்பும் உள்ளது); மேலும்
ஒப்பிடத்தக்க பண்பும் உள்ளது. ஆனால் எதிர்சமச்சீர்ப் பண்பு இல்லாததால் இது
பகுதி வரிசையல்ல. எடுத்துக்காட்டாக, $01
\preccurlyeq 10$, $10 \preccurlyeq 01$ ஆகிய இரண்டும் உண்மை; ஆனால் $01 \neq 10$.
\end{ex}

% SEGMENT OLP-0016-S07
\begin{ex}
கணங்களைக் கொண்ட ஒரு கணத்தின் மீதான $\subseteq$ தொடர்பு ஒரு முக்கியமான
பகுதி வரிசை. பொதுவாக இது நேரியல் வரிசையல்ல. ஏனெனில் $a \neq b$
என வைத்து $\Pow{\{a, b\}} = \{\emptyset, \{a\}, \{b\}, \{a,b\}\}$
ஐக் கருதினால், $\{a\} \nsubseteq \{b\}$, $\{a\} \neq \{b\}$,
$\{b\} \nsubseteq \{a\}$ ஆகியவற்றைக் காணலாம்.
\end{ex}

% SEGMENT OLP-0016-S08
\begin{ex}
\emph{மீதியின்றி வகுபடும்} தொடர்பு, நேரியல் வரிசையாக இல்லாத ஒரு பகுதி
வரிசையைத் தருகிறது. முழுக்களான $n$, $m$ க்கு, $n$ என்பது $m$ ஐ
மீதியின்றி வகுக்கிறது என்பதைக் குறிக்க $n \mid m$ என்று எழுதுகிறோம்.
அதாவது $m = kn$ என அமையும் ஒரு முழு $k$ இருக்கும்போது,
அப்போதும் அப்போது மட்டுமே, இவ்வாறு எழுதுகிறோம். $\Nat$ மீது இது
பகுதி வரிசை; ஆனால் நேரியல் வரிசையல்ல: எடுத்துக்காட்டாக, $2 \nmid 3$,
அதேபோல $3 \nmid 2$. $\Int$ மீதான தொடர்பாகக் கருதும்போது, வகுபடுதல்
ஒரு முன்வரிசை மட்டுமே. ஏனெனில் அதற்கு எதிர்சமச்சீர்ப் பண்பு இல்லை:
$1 \mid -1$, $-1 \mid 1$ ஆகிய இரண்டும் உண்மை; ஆனால் $1 \neq -1$.
\end{ex}

% SEGMENT OLP-0016-S09
\begin{ex}
தொடர்முறைகளின் கணமான $A^*$ மீதான \emph{நீட்டிப்புத்} தொடர்பு
பின்வருமாறு அமைகிறது: $s = \emptyseq$ (வெற்றுத் தொடர்முறை), $s = s'$,
அல்லது $s = \tuple{s_1, \dots, s_n}$ மற்றும் $s' =
\tuple{s_1, \dots, s_n, s_{n+1}, \dots, s_m}$ ஆகிய நிலைகளில் ஒன்று
இருக்கும்போது, அப்போதும் அப்போது மட்டுமே, $s \sqsubseteq s'$.
$s \sqsubseteq s'$ எனில், $s$ என்பது $s'$ இன் \emph{தொடக்கத் துண்டு}
என்றும் கூறுகிறோம். $A^*$ மீதான நீட்டிப்புத் தொடர்பு ஒரு பகுதி வரிசை;
ஆனால் நேரியல் வரிசையல்ல. எ.கா., $a \neq b$ எனில்,
$ab \not\sqsubseteq ba$, $ba \not\sqsubseteq ab$.
\end{ex}

% SEGMENT OLP-0016-S10
\begin{defn}[கண்டிப்பான வரிசை]
எங்கும் தற்சுட்டற்ற பண்பு, முழுச் சமச்சீரின்மை, கடப்புப் பண்பு ஆகியவற்றைக்
கொண்ட ஒரு தொடர்பு \emph{கண்டிப்பான வரிசை} ஆகும்.
\end{defn}

% SEGMENT OLP-0016-S11
\begin{defn}[கண்டிப்பான நேரியல் வரிசை]\ollabel{def:strictlinearorder}
ஒப்பிடத்தக்க பண்பையும் கொண்ட ஒரு கண்டிப்பான வரிசை,
\emph{கண்டிப்பான முழு வரிசை} அல்லது \emph{கண்டிப்பான நேரியல் வரிசை}
எனப்படும்.
\end{defn}

% SEGMENT OLP-0016-S12
\begin{ex}
$\le$ என்பது கண்டிப்பான நேரியல் வரிசையான $<$ க்கு உரிய நேரியல் வரிசை.
$\subseteq$ என்பது கண்டிப்பான வரிசையான $\subsetneq$ க்கு உரிய பகுதி வரிசை.
\end{ex}

% SEGMENT OLP-0016-S13
$A$ மீதான எந்தக் கண்டிப்பான வரிசை $R$ ஐயும், மூலைவிட்டமான $\Id{A}$
ஐச் சேர்ப்பதன் மூலம், அதாவது $\tuple{x, x}$ என்ற அனைத்து
ஜோடிகளையும் சேர்ப்பதன் மூலம், பகுதி வரிசையாக மாற்றலாம்.
(இது $R$ இன் \emph{தற்சுட்டு அடைவு} எனப்படுகிறது.)
மறுதலையாக, ஒரு பகுதி வரிசையிலிருந்து $\Id{A}$ ஐ நீக்கினால் ஒரு
கண்டிப்பான வரிசை கிடைக்கும். அடுத்து வரும் இரு முடிவுகள் இதைத்
துல்லியமாகக் கூறுகின்றன.

% SEGMENT OLP-0016-S14
\begin{prop}\ollabel{prop:stricttopartial}
$R$ என்பது $A$ மீதான ஒரு கண்டிப்பான வரிசை எனில்,
$R^+ = R \cup \Id{A}$ ஒரு பகுதி வரிசை. மேலும், $R$ ஒரு கண்டிப்பான
நேரியல் வரிசை எனில், $R^+$ ஒரு நேரியல் வரிசை.
\end{prop}

% SEGMENT OLP-0016-S15
\begin{proof}
$R$ ஒரு கண்டிப்பான வரிசை எனக் கொள்வோம்; அதாவது $R \subseteq A^2$,
மேலும் $R$ எங்கும் தற்சுட்டற்ற பண்பு, முழுச் சமச்சீரின்மை, கடப்புப் பண்பு
ஆகியவற்றைக் கொண்டுள்ளது. $R^+ = R \cup \Id{A}$ என வைப்போம்.
$R^+$ தற்சுட்டுப் பண்பு, எதிர்சமச்சீர்ப் பண்பு, கடப்புப் பண்பு
ஆகியவற்றைக் கொண்டுள்ளது எனக் காட்ட வேண்டும்.

எல்லா $x \in A$ க்கும் $\tuple{x, x} \in \Id{A} \subseteq
R^+$ என்பதால், $R^+$ தற்சுட்டுப் பண்பைக் கொண்டுள்ளது என்பது தெளிவு.

% SEGMENT OLP-0016-S16
$R^+$ எதிர்சமச்சீர்ப் பண்பைக் கொண்டுள்ளது எனக் காட்ட, முரண் மூலம்
நிறுவும் பொருட்டு $R^+xy$, $R^+yx$, ஆனால் $x \neq y$ எனக் கொள்வோம்.
$\tuple{x,y} \in R \cup \Id{A}$, ஆனால்
$\tuple{x, y} \notin \Id{A}$ என்பதால்,
$\tuple{x, y} \in R$, அதாவது $Rxy$, இருக்க வேண்டும். அதேபோல
$Ryx$. ஆனால் இது $R$ முழுச் சமச்சீரின்மையைக் கொண்டுள்ளது என்ற
அனுமானத்துடன் முரண்படுகிறது.

% SEGMENT OLP-0016-S17
கடப்புப் பண்பை நிறுவ, $R^+xy$, $R^+yz$ எனக் கொள்வோம்.
$\tuple{x, y} \in R$, $\tuple{y,z} \in R$ ஆகிய இரண்டும் உண்மையெனில்,
$R$ கடப்புப் பண்பைக் கொண்டுள்ளதால் $\tuple{x, z} \in
R$. இல்லாவிட்டால், $\tuple{x, y} \in
\Id{A}$, அதாவது $x = y$, அல்லது $\tuple{y, z} \in \Id{A}$,
அதாவது $y = z$. முதல் நிலையில், அனுமானத்தின்படி $R^+yz$; மேலும்
$x = y$ என்பதால் $R^+xz$. இரண்டாம் நிலையிலும் இதேபோலக் காட்டலாம்.
இரு நிலைகளிலும் $R^+xz$; ஆகவே $R^+$ கடப்புப் பண்பையும் கொண்டுள்ளது.

% SEGMENT OLP-0016-S18
“மேலும்” என்ற பகுதிக்காக, $R$ ஒப்பிடத்தக்க பண்பையும் கொண்டுள்ளது
எனக் கொள்வோம். அப்போது எல்லா $x \neq y$ க்கும் $Rxy$ அல்லது $Ryx$;
அதாவது $\tuple{x, y} \in R$ அல்லது $\tuple{y, x} \in R$.
$R \subseteq R^+$ என்பதால், $R^+$ க்கும் இது உண்மையாகவே உள்ளது.
எனவே $R^+$ ஒப்பிடத்தக்க பண்பையும் கொண்டுள்ளது.
\end{proof}

% SEGMENT OLP-0016-S19
\begin{prop}\ollabel{prop:partialtostrict}
$R$ என்பது $A$ மீதான ஒரு பகுதி வரிசை எனில்,
$R^- = R \setminus \Id{A}$ ஒரு கண்டிப்பான வரிசை. மேலும், $R$ ஒரு
நேரியல் வரிசை எனில், $R^-$ ஒரு கண்டிப்பான நேரியல் வரிசை.
\end{prop}

% SEGMENT OLP-0016-S20
\begin{proof}
இது ஒரு பயிற்சியாக விடப்படுகிறது.
\end{proof}

\begin{prob}
\olref[sfr][rel][ord]{prop:partialtostrict} க்கு ஒரு நிறுவல் தருக.
\end{prob}

% SEGMENT OLP-0016-S21
பின்வரும் எளிய முடிவு, கண்டிப்பான நேரியல் வரிசைகள் உறுப்புசார் சமத்துவம்
போன்ற ஒரு பண்பை நிறைவு செய்கின்றன என்பதை நிறுவுகிறது:

\begin{prop}\ollabel{prop:extensionality-strictlinearorders}
$<$ என்பது $A$ மீதான ஒரு கண்டிப்பான நேரியல் வரிசை எனில்:
\[
  (\forall a, b \in A)((\forall x \in A)(x < a \liff x < b) \lif a = b).
\]
\end{prop}

% SEGMENT OLP-0016-S22
\begin{proof}
$(\forall x \in A)(x < a \liff x < b)$ எனக் கொள்வோம். $a < b$ எனில்,
$a < a$ ஆகிறது; இது $<$ எங்கும் தற்சுட்டற்றது என்பதுடன் முரண்படுகிறது.
எனவே $a \nless b$. முற்றிலும் இதேபோல $b \nless a$.
$<$ ஒப்பிடத்தக்க பண்பைக் கொண்டுள்ளதால் $a = b$.
\end{proof}

\end{document}
